\subsection*{用 \texttt{swap} 做双缓冲}

给 K 循环准备两个逻辑缓冲区：计算使用缓冲区 0 时，DMA 将下一个分块装入缓冲区 1。计算完成后交换句柄，原来的“下一个”变成“当前”，释放的槽位供下一次装入使用。

在 Croqtile 中，这一过程由 \co@dma.copy.async@（非阻塞拷贝）、\co@dma.any@（占位期值）、\co@swap@（交换期值句柄）和一个三阶段循环来表达：

\begin{lstlisting}[language=croqtile,caption={用 swap 做双缓冲}]
__co__ auto matmul(s32 [M, K] lhs, s32 [K, N] rhs) {
  s32 [lhs.span(0), rhs.span(1)] output;
  parallel {px, py} by [8, 16] : block
    parallel {qx, qy} by [16, 16] : thread {
    output.at(px#qx, py#qy) = 0;
    with tile_k in 16 {
      // Prologue: start loading tile 0
      lf0 = dma.copy.async lhs.chunkat(px, tile_k) => shared;
      rf0 = dma.copy.async rhs.chunkat(tile_k, py) => shared;
      // Placeholder futures for buffer 1
      lf1 = dma.any;
      rf1 = dma.any;
      // Steady state: load next tile while computing on current
      foreach tile_k(1:) {
        lf1 = dma.copy.async lhs.chunkat(px, tile_k) => shared;
        rf1 = dma.copy.async rhs.chunkat(tile_k, py) => shared;
        wait lf0, rf0;
        foreach k in [256 / #tile_k]
          output.at(px#qx, py#qy) += lf0.data.at(qx, k) * rf0.data.at(k, qy);
        swap(lf0, lf1);
        swap(rf0, rf1);
      }
      // Epilogue: compute on the last loaded tile
      foreach k in [256 / #tile_k]
        output.at(px#qx, py#qy) += lf0.data.at(qx, k) * rf0.data.at(k, qy);
    }
  }
  return output;
}
\end{lstlisting}
